\section{Formalization Map}
\label{sec:formalization-map}

This section maps the informal proof to the current Lean project structure.

\begin{center}
\begin{tabular}{p{0.34\linewidth}p{0.56\linewidth}}
\textbf{Challenge declaration} & \textbf{Informal source} \\
\hline
\code{genus} &
Definition~\ref{def:analytic-genus}; finite-dimensionality from
Input~\ref{input:finite-dimensionality}. \\
\code{genus_eq_zero_iff_homeo} &
Theorem~\ref{thm:genus-zero-classification}; uses Riemann-Roch and
degree-one map classification. \\
\code{Jacobian} and basic instances &
Definition~\ref{def:analytic-jacobian}; period lattice theorem plus complex
torus package. \\
\code{Jacobian.ofCurve} &
Definition~\ref{def:abel-jacobi}; path independence modulo periods. \\
\code{ofCurve_contMDiff} &
Lemma~\ref{lem:aj-holomorphic}. \\
\code{ofCurve_self} &
Lemma~\ref{lem:aj-self}. \\
\code{ofCurve_inj} &
Theorem~\ref{thm:abel-jacobi-injective}; Abel's theorem. \\
\code{pushforward}, \code{pullback} &
Lemma~\ref{lem:push-pull-descend}; trace and pullback preserving periods. \\
\code{pushforward_contMDiff} and \code{pullback_contMDiff} &
Theorem~\ref{thm:push-pull-functoriality}; descended complex-linear maps are
holomorphic. \\
\code{pushforward_id_apply} and \code{pullback_id_apply} &
Theorem~\ref{thm:push-pull-functoriality}; identity laws for pullback and
trace. \\
\code{pushforward_comp_apply} and \code{pullback_comp_apply} &
Theorem~\ref{thm:push-pull-functoriality}; composition laws for pullback and
trace. \\
\code{ContMDiff.degree} &
Definition~\ref{def:degree-trace}. \\
\code{pushforward_pullback} &
Theorem~\ref{thm:pushforward-pullback}; trace-pullback identity. \\
\end{tabular}
\end{center}

The current repo already has production modules for many algebraic and
topological parts of this map. The largest classical gaps are:

\begin{itemize}
\item Riemann-Roch and divisor theory for compact Riemann surfaces.
\item Abel's theorem in point-separation form.
\item Riemann bilinear relations and full-rank period lattice theorem.
\item Branched-cover degree theory and trace of holomorphic differentials.
\item Analytic genus versus topological genus comparison.
\end{itemize}
